{{دیوان
| عنوان = مذہب (تعلیمِ پیرِ فلسفۂ مغربی ہے یہ)
| ماخذ = ویکی ماخذ
| ماخذ_لنک = https://ur.wikisource.org/wiki/%D9%85%D8%B0%DB%81%D8%A8_%28%D8%AA%D8%B9%D9%84%DB%8C%D9%85%D9%90_%D9%BE%DB%8C%D8%B1%D9%90_%D9%81%D9%84%D8%B3%D9%81%DB%82_%D9%85%D8%BA%D8%B1%D8%A8%DB%8C_%DB%81%DB%92_%DB%8C%DB%81%29
| مقابلہ = اقبال سائبر لائبریری (اقبال اکادمی پاکستان)
| مقابلہ_لنک = https://www.iqbalcyberlibrary.net
}}
=== تضمین بر شعرِ میرزا بیدلؔ ===

<poem>
تعلیمِ پیرِ فلسفۂ مغربی ہے یہ
ناداں ہیں جن کو ہستیِ غائب کی ہے تلاش

پیکر اگر نظر سے نہ ہو آشنا تو کیا
ہے شیخ بھی مثالِ بَرہمن صنَم تراش

محسوس پر بِنا ہے علومِ جدید کی
اس دَور میں ہے شیشہ عقائد کا پاش پاش

مذہب ہے جس کا نام، وہ ہے اک جنُونِ خام
ہے جس سے آدمی کے تخیّل کو انتعاش

کہتا مگر ہے فلسفۂ زندگی کچھ اور
مجھ پر کِیا یہ مُرشدِ کامل نے راز فاش

“باہر کمال اند کے آشفتگی خوش است
ہر چند عقلِ کُل شدہ ای بے جُنوں مباش”
</poem>
